\documentclass[a4paper]{article}

\usepackage{graphicx}
\usepackage{amsmath,amssymb}
\usepackage{minted}
\usepackage{multirow}
\usepackage{float}
\usepackage{todonotes}
\usepackage{forest}
\usepackage{xstring}
\usepackage{xcolor}
\usepackage[most]{tcolorbox}
\usepackage{xparse}
\usepackage{natbib}
\usepackage{algorithm}
\usepackage{algpseudocode}

\usetikzlibrary{graphs,graphdrawing}
\usegdlibrary{layered}

% for external graphviz
\input{/home/donkarlo/Dropbox/repo/nd_ai_project/src/nd_ai/knoweldge_representation/language/formal/computer/programming/tex/latex/code_clips/external_graphviz.tex}

% for tags
\input{/home/donkarlo/Dropbox/repo/nd_ai_project/src/nd_ai/knoweldge_representation/language/formal/computer/programming/tex/latex/parts/control_sequence/kind/semtag/semtag.tex}

% for links
\input{/home/donkarlo/Dropbox/repo/nd_ai_project/src/nd_ai/knoweldge_representation/language/formal/computer/programming/tex/latex/code_clips/seehere.tex}

\title{Choosing number of classes}
\date{}
\author{Mohammad Rahmani}

\begin{document}
    \maketitle
    
    
    \section{Silhouette score}
        The silhouette score \cite{rousseeuw-1987-silhouettes-a-graphical-aid-to-the-interpretation-and-validation-of-cluster-analysis} measures how well each sample belongs to its assigned cluster in comparison with the nearest alternative cluster. For a sample \(i\), let \(a_i\) denote the mean distance between \(i\) and the other samples in its own cluster, and let \(b_i\) denote the smallest mean distance between \(i\) and samples in another cluster. The silhouette value is defined as
        
        \[
            s_i
            =
            \frac{b_i - a_i}{\max(a_i, b_i)}.
        \]
        
        For a clustering solution with \(k\) clusters, the average silhouette score is
        
        \[
            S(k)
            =
            \frac{1}{n}
            \sum_{i=1}^{n}
            s_i.
        \]
        
        Since the ordinary silhouette criterion may favor a larger number of clusters when the improvement is small, a penalty term can be added to discourage unnecessary cluster splitting:
        
        \[
            S_{\mathrm{pen}}(k)
            =
            S(k)
            -
            \lambda
            \frac{k - 1}{k_{\max} - 1},
        \]
        
        where \(\lambda\) controls the strength of the penalty, and \(k_{\max}\) is the maximum tested number of clusters. The final number of clusters is selected as
        
        \[
            k^{\ast}
            =
            \arg\max_{k}
            S_{\mathrm{pen}}(k).
        \]
        
        In this analysis, the penalized silhouette criterion selected
        
        \[
            k^{\ast} = 5.
        \]
        
        \subsection{penalized average silhouette width (ASW)}
            \cite{deng-2024-on-estimation-and-order-selection-for-multivariate-extremes-via-clustering} propose a penalized average silhouette width (ASW) criterion for selecting the number of clusters in multivariate extreme-value analysis.
            The penalty is designed to reduce spurious cluster selection caused by tiny isolated clusters or by splitting one true center into several nearby centers.
            The method selects the cluster order by inspecting penalized ASW curves over small penalty values, rather than by relying only on the ordinary silhouette maximum.
            \begin{algorithm}[H]
            \caption{Penalized ASW Order Selection}
            \begin{algorithmic}[1]
                \State Choose candidate orders \(k \in \{1,\dots,k_{\max}\}\) and small penalty values \(t \in T\).
                \For{each \(k\)}
                    \State Cluster the sample \(W\) and obtain centers \(A_k^\ast=\{a_1^\ast,\dots,a_k^\ast\}\) and clusters \(C_1,\dots,C_k\).
                    \State For each \(w \in W\), compute \(a(w)=D(w,A_k^\ast)\) and \(b(w)=D(w,A_k^\ast \setminus \{a_i^\ast\})\).
                    \State Compute \(\bar{s}(k)=1-\frac{1}{|W|}\sum_{w\in W}\frac{a(w)}{b(w)}\).
                    \State Compute the penalty factor \(p_t(k)=\left(\min_i \frac{|C_i|}{|W|/k}\right)^t\left(\min_{i<j}D(a_i^\ast,a_j^\ast)\right)^t\).
                    \State Compute \(S_t(k)=p_t(k)-\frac{1}{|W|}\sum_{w\in W}\frac{a(w)}{b(w)}\).
                \EndFor
                \State Select \(k^\ast\) from the maximum or turning point of \(S_t(k)\) over small \(t\) values.
            \end{algorithmic}
            \end{algorithm}
            
            \bibliographystyle{plainnat}
            \bibliography{/home/donkarlo/Dropbox/repo/nd_shared_research/bibliography/bibliography.bib}
\end{document}